% from main.tex
%%%%%%%%%%  scan idx 279  |  folio 266  %%%%%%%%%%
Si \(E\)
vérifie la condition (5.1 a)) pour l'entier \(p\), \(E^{\vee}\) vérifie la même condition pour l'entier \(-p-1\) en vertu de (5.2 (iii)), et la flèche \((*)\) est un isomorphisme, d'après la première assertion de (5.3.1). Ceci achève la démonstration de (5.3.1), et par suite de (5.3).

\textbf{Remarque 5.3.2.} La proposition (5.3) s'applique en particulier lorsque \(E\) est un complexe à degrés bornés de \(\mathcal O_X\)-Modules quasi-cohérents plats et de présentation finie rel. à \(f\) : en effet, \(E\) est alors d'amplitude parfaite finie rel. à \(f\) en vertu de (4.3.2).

